\documentclass[journal]{IEEEtran}
\usepackage{iftex}
\ifPDFTeX
\usepackage[utf8]{inputenc}
\usepackage[T1]{fontenc}
\fi
\usepackage{amsmath,amssymb,amsthm,bm}
\usepackage{newtxtext,newtxmath}
\usepackage{booktabs,array,adjustbox}
\usepackage{balance}
\usepackage{siunitx}
\usepackage[final]{microtype}
\usepackage[breaklinks=true,hidelinks]{hyperref}
\interdisplaylinepenalty=2500
\allowdisplaybreaks
\emergencystretch=2em
\newtheorem{lemma}{Lemma}
\newtheorem{theorem}{Theorem}
\newtheorem{remark}{Remark}
\newcommand{\R}{\mathbb{R}}
\newcommand{\norm}[1]{\left\lVert#1\right\rVert}
\newcommand{\Id}{I}
\newcommand{\fitcol}[1]{\adjustbox{max width=\columnwidth}{\ensuremath{\displaystyle #1}}}
\AtBeginEnvironment{equation}{\small}
\AtBeginEnvironment{align}{\small}

\title{Stability Study of an Adaptive OU-Regularized Quaternion MEKF for Marine Inertial Wave Estimation}
\author{Mikhail S. Grushinskiy}

\begin{document}
\maketitle

\begin{abstract}
This study formulates a deterministic stability architecture for the shipping
OU--III marine estimator. The physical domain is defined by three simultaneous
contracts: bounded marine motion with an all-time bounded wave-displacement
primitive, bounded and rate-bounded total IMU residual bias, and recurring
accepted informative magnetic observations. The estimator bias prior is kept
distinct from physical bias truth. One prediction relation covers held and
active accelerometer-bias modes, while correction and estimate projection are
retained as separate implemented operations. The analysis uses one path from
construction through finite capture, held-bias Live operation, bias-state
release, and magnetically informed active-bias operation. An exact unipotent
heading/axial gyro-bias block establishes the need for continuing heading
information, and an identity block on same-history held-bias increments excludes strict
full-state incremental contraction on that leg. A conditional complement
inequality retains held bias and finite-error remainders as inputs. End-to-end
regional practical stability remains open until finite capture, finite-error
dissipativity, retention, and implementation arithmetic are certified.
\end{abstract}

\section{Theorem scope}
Let one physical execution be denoted by $h$. The principal domain is
\begin{equation}
 h\in\mathcal M_{\rm marine}\cap\mathcal B_{\rm IMU}\cap\mathcal S_{\rm mag}.
 \label{eq:intersection}
\end{equation}
The intersection is a same-history requirement. The execution determines
vessel translation, attitude, IMU measurements, physical biases, magnetic
measurements, frontend and tuning state, covariance, event acceptance,
scheduler state, and estimator variables. The proof boundary does not alter
the deployed estimator.

\section{Marine motion}
Let $p(t)\in\R^3$ be wave-induced displacement about a local equilibrium,
$v=\dot p$, $a=\dot v$, with the same history supplying
$R(t)\in{\rm SO}(3)$ and body rate $\omega(t)$. Here $R$ maps
body to world and satisfies $\dot R=R[\omega]_\times$ almost everywhere. Require
\begin{align}
 \norm{p(t)}&\le P_{\max},& \norm{v(t)}&\le V_{\max},\nonumber\\
 \norm{a(t)}&\le A_{\max},& \norm{\omega(t)}&\le\Omega_{\max}.
 \label{eq:pointwise}
\end{align}
Pointwise bounds are insufficient. Admission additionally requires
\begin{equation}
 \boxed{\norm{\int_{t_1}^{t_2}p(\tau)\,d\tau}\le P_{\rm AC},
 \qquad t_2\ge t_1 .}
 \label{eq:pac}
\end{equation}
Equivalently, $\dot q=p$ has uniformly bounded potential differences. The
same potential continues through successive certified superwords; it is not
restarted at word boundaries. A one-time proof-coordinate origin at tail entry
is permissible only as a coordinate transformation of that same continuation.
With the origin fixed at capture, put $S_{\rm true}(t)=q(t)-q(T_c)$.
Consequently $\norm{S_{\rm true}}\le P_{\rm AC}$. The implemented
zero-integral innovation is
\begin{equation}
 r_S=-\hat S=e_S-S_{\rm true},\qquad e_S=S_{\rm true}-\hat S.
 \label{eq:physical-integral-innovation}
\end{equation}
The physical term $-S_{\rm true}$ is therefore retained through every actual
integral correction. Neither $\hat S$ nor the physical potential is reset at
certification or at later superword boundaries.


\begin{lemma}[Permanent displacement DC is inadmissible]
If $p(t)\equiv p_0\ne0$, $v(t)=0$, and $a(t)=0$ indefinitely, then
\eqref{eq:pac} fails.
\end{lemma}
\begin{proof}
On an interval of length $T$,
$\norm{\int p\,dt}=T\norm{p_0}\to\infty$.
\end{proof}
Quiet water $p\equiv0$ remains admissible. Finite approximately constant
segments are not excluded. Bounded samples, a finite mean, a spectrum, or
finite-window increments do not prove \eqref{eq:pac}.

Physical CoG translation may be written
\begin{equation}
 p_{\rm CoG}=p_{\rm eq}+p .
\end{equation}
Global origin, current, propulsion, leeway, and secular reference translation
may be represented by $p_{\rm eq}$. If $\ddot p_{\rm eq}\ne0$, however, the
associated specific force remains an implemented model term or a declared
same-history disturbance.

\section{IMU bias}
The total residual accelerometer and gyro biases after the calibration
actually performed are locally absolutely continuous and satisfy
\begin{align}
 \norm{b_a(t)}&\le B_a,& \norm{\dot b_a(t)}&\le D_a,\nonumber\\
 \norm{b_g(t)}&\le B_g,& \norm{\dot b_g(t)}&\le D_g
 \quad\text{a.e.}
 \label{eq:bias-cont}
\end{align}
The constants are qualified separately. No decomposition of physical truth is
imposed: the aggregate residual may simultaneously include turn-on offset,
thermal drift, strain, hysteresis, aging, correlated drift, and other
calibration residuals. The sampled physical histories obey
\begin{align}
 b_{a,k+1}&=b_{a,k}+w_{a,k},&
 \norm{w_{a,k}}&\le D_a\Delta t_k,\nonumber\\
 b_{g,k+1}&=b_{g,k}+w_{g,k},&
 \norm{w_{g,k}}&\le D_g\Delta t_k .
 \label{eq:bias-discrete}
\end{align}
Thus the successor is constrained by its predecessor rather than selected
independently from an amplitude box.

Let $e_b=b-\hat b$. For the accelerometer-bias estimate, define the implemented
prediction factor
\begin{equation}
 \phi_{e,k}=
 \begin{cases}
 1,&\text{H18/held prediction},\\
 \phi_{{\rm OU},k},&\text{A21/active prediction}.
 \end{cases}
 \label{eq:phi-mode}
\end{equation}
With physical truth $b_{k+1}=b_k+w_k$, both modes obey the single relation
\begin{equation}
 \boxed{
 e_{b,k+1}^{-}=
 \phi_{e,k}e_{b,k}^{+}+
 (1-\phi_{e,k})b_k+w_k .}
 \label{eq:bias-prediction}
\end{equation}
Equation~\eqref{eq:bias-prediction} is an estimator-mode relation, not a model
imposed on physical truth.
For $z_b=(e_b,b)$ the same identity has the joint form
\begin{equation}
 \begin{bmatrix}e_{b,k+1}^{-}\\b_{k+1}\end{bmatrix}
 =\begin{bmatrix}\phi_{e,k}I&(1-\phi_{e,k})I\\0&I\end{bmatrix}
 \begin{bmatrix}e_{b,k}^{+}\\b_k\end{bmatrix}
 +\begin{bmatrix}I\\I\end{bmatrix}w_k.
 \label{eq:bias-joint}
\end{equation}
The two rows share one physical increment. Neither independent disturbance
columns nor an estimator-induced reset of $b$ preserves this relation.


Measurement correction is carried separately. If the implemented correction
changes the estimate by $\Delta\hat b_k$, then
\begin{equation}
 e_{b,k}^{\rm corr}=e_{b,k}^{-}-\Delta\hat b_k .
 \label{eq:bias-correction}
\end{equation}
For the shipping accelerometer-bias projection,
\begin{align}
 \hat b_{a,k}^{+}&=\operatorname{Proj}_{R_b}
                      (\hat b_{a,k}^{\rm corr}),\nonumber\\
 e_{b_a,k}^{+}&=b_{a,k}-
 \operatorname{Proj}_{R_b}(\hat b_{a,k}^{\rm corr}).
 \label{eq:bias-projection}
\end{align}
For a finite completed injection with $R_b>0$, write
$d_{b,k}=\hat b_{a,k}^{\rm corr}-\operatorname{Proj}_{R_b}
(\hat b_{a,k}^{\rm corr})$. Then
\begin{equation}
 e_{b_a,k}^{+}=e_{b_a,k}^{\rm corr}+d_{b,k}.
 \label{eq:bias-defect}
\end{equation}
\begin{lemma}[Estimate-projection bounds]
For $\norm{b_a}\le B_a$ and an applied positive-radius projection,
$\norm{e_{b_a}^{+}}\le B_a+R_b$. If additionally $B_a\le R_b$, then
\begin{equation}
 \norm{e_{b_a}^{+}}^2+\norm{d_b}^2
 \le\norm{e_{b_a}^{\rm corr}}^2.
 \label{eq:bias-projection-sector}
\end{equation}
\end{lemma}
\begin{proof}
The first bound is the triangle inequality. The projection variational
inequality gives $(e_{b_a}^{+})^{\mathsf T}d_b\le0$ when $b_a$ lies in
the ball. Expanding $e_{b_a}^{\rm corr}=e_{b_a}^{+}-d_b$ proves the second.
\end{proof}
This Euclidean component inequality is not contraction in a full-state
covariance-weighted metric. Cross terms and all covariance operations remain
in the finite-error argument. A nonpositive configured radius disables the
shipping projection. Moreover, an invalid attitude injection returns before
projection; excluding or enclosing that branch is an arithmetic/domain
obligation. The H18--A21 release preserves nominal bias states but applies
its covariance floor. Any frame change must transport both physical bias and
its estimate before the next prediction.

Physical bias is unchanged by prediction mode, correction, hold/release, or
projection. The shipping gyro-bias mean predictor is identity, giving
$e_{b_g,k+1}^{-}=e_{b_g,k}^{+}+w_{g,k}$ before its actual correction.

\section{Magnetic service}
For a certified service window rooted at $s$, let $\Phi_{k\leftarrow s}$ be
the ordered differential of the complete preceding same-history shipping
evolution. Let $H_{m,k}$ be the literal magnetic sensitivity and let
$S_{m,k}^{\rm act}$ be the actual innovation covariance presented to the
shipping update's factorization for a correction that was actually applied.
Choose any whitening factor $W_k$ satisfying
$W_k^\mathsf{T}W_k=(S_{m,k}^{\rm act})^{-1}$. With $E_{hb}$ injecting
normalized heading and axial gyro-bias coordinates at the window root, define
\begin{equation}
 G_k=W_kH_{m,k}\Phi_{k\leftarrow s}E_{hb}.
 \label{eq:G}
\end{equation}
Equivalently, if $S_{m,k}^{\rm act}=L_kL_k^\mathsf{T}$, one may take
$W_k=L_k^{-1}$. This square-root choice is algebraic only; the information
matrix is built from the same $S_{m,k}^{\rm act}$ that the shipping update
actually factors. Transport may be rectangular at an error-coordinate
change: $H_{m,k}$ uses destination coordinates and $E_{hb}$ uses the
original root coordinates. This algebra does not certify a release map or
a finite-error information margin. For
$T_M<\infty$ and $\mu_M>0$, every certified tail interval must satisfy
\begin{equation}
 \boxed{\sum_{k\in\mathcal M[s,s+T_M]}G_k^\mathsf{T}G_k
 \succeq\mu_M I_2.}
 \label{eq:mag-service}
\end{equation}
Only finite, unsaturated, gauged measurements whose shipping correction was
actually applied belong to $\mathcal M$. Attempted callbacks, due events,
packets, rejected or invalid measurements, and saturation do not contribute.
A maximum event gap alone is insufficient because frequent collinear
sensitivities may remain rank deficient.

\section{Single shipping proof path}
The theorem uses
\begin{equation}
\begin{aligned}
 &\text{construction}\to\text{capture}\to\text{informed H18}\\
 &\quad\to\text{H18--A21 release}\to\text{informed A21}\\
 &\quad\to\text{practical stability}.
\end{aligned}
\end{equation}
The physical bias history is continuous through all modes. The capture time may
depend on execution and initial state,
\begin{equation}
 \boxed{T_c=T_c(h,x_0)<\infty.}
\end{equation}
No common finite startup deadline is assumed. Live intervals preceding the
required magnetic information remain part of the finite pre-certified prefix.
At tail entry the proof inherits, without reseeding, the actual estimator
state, full covariance, bias estimates, physical bias histories, frontend and
tuner state, committed parameters, magnetic reference state, scheduler,
clocks, and physical source continuation.

\section{Necessity of heading information}
On an interval without informative heading observations, heading error and
axial gyro-bias error contain
\begin{equation}
 \begin{bmatrix}\theta_\parallel^+\\e_{b_g,\parallel}^+\end{bmatrix}
 =
 \begin{bmatrix}1&T\\0&1\end{bmatrix}
 \begin{bmatrix}\theta_\parallel\\e_{b_g,\parallel}\end{bmatrix}.
 \label{eq:unipotent}
\end{equation}
Both eigenvalues are one and repeated application gives linear heading growth
when $e_{b_g,\parallel}\ne0$. Hence strict coercive full-state contraction is
not generally available during indefinite loss of informative heading service.
This is a necessity result, not a second stability theorem.

\section{Non-contraction of the held bias}
On a held segment with a feasible bias estimate and no projection-changing,
frame-change or relock event, the bias estimate is constant. Its gain rows
vanish and its decoupled covariance block is unchanged. Absolute bias error
still satisfies $e_b^+=e_b+w$; physical increments cancel only between
executions on the same physical history. The incremental differential is
\begin{equation}
 \Psi=\begin{bmatrix}A&B\\0&\Id\end{bmatrix}.
 \label{eq:held-bias-block}
\end{equation}
This is not a global linear formula for the nonlinear finite-error map.

\begin{lemma}[Held bias blocks strict contraction]
Under \eqref{eq:held-bias-block} with $P_{bb}$ unchanged and the bias
cross-covariances zero, the incremental storage satisfies
\begin{equation}
 \max_{e\ne0}\frac{V_{\rm end}(\Psi e)}{V_{\rm root}(e)}\ge1 .
 \label{eq:held-bias-obstruction}
\end{equation}
\end{lemma}
\begin{proof}
Take $e=(0,e_b)$. Block diagonality gives
$V_{\rm root}(e)=e_b^\mathsf{T}P_{bb}^{-1}e_b$ and
$V_{\rm end}(\Psi e)=(Be_b)^\mathsf{T}P_{oo,\rm end}^{-1}(Be_b)
+e_b^\mathsf{T}P_{bb}^{-1}e_b$, since $P_{bb}$ is unchanged. The first term is
nonnegative.
\end{proof}
Since $\Psi$ is block triangular with an identity block it also carries the
eigenvalue one, so no time-invariant quadratic storage attains $\rho<1$ either.
The obstruction holds whenever the stated segment premises hold. Hard events
and release require their own retention analysis.

This is a necessity result, not an instability claim. Under the completed
projection premises, $\norm{e_b}\le B_a+R_b$. The H18 leg therefore retains
held bias as a bounded input on the complement; the actual release and A21
analysis follow on the same proof path.

\section{Complement storage and finite-error supply}
For complement error $x$, retain the exact relation
\begin{equation}
 x^+=Ax+Be_b+r,\qquad V=x^{\mathsf T}Wx,
 \label{eq:complement-map}
\end{equation}
where $W=P_{oo}^{-1}$. For a chosen local surrogate $A,B$, this relation
defines $r$; it does not bound it. Reference forcing, physical bias increments,
nonlinear terms, auxiliary-state dependence, model mismatch and arithmetic
remainders must remain represented in the supply.

\begin{lemma}[Conditional complement supply]
Suppose $A^{\mathsf T}W_+A\preceq\alpha W$ on a retained same-history domain,
with $0\le\alpha<\rho<1$. Then
\begin{equation}
 V_+\le\rho V+\frac{\rho}{\rho-\alpha}
          \norm{Be_b+r}_{W_+}^2.
 \label{eq:held-supply}
\end{equation}
\end{lemma}
\begin{proof}
Bound $\norm{Ax}_{W_+}^2$ by $\alpha V$ and apply weighted Young's inequality
to the cross term with weight $(\rho-\alpha)/\alpha$.
The case $\alpha=0$ follows directly.
\end{proof}

For the isolated linear bias input, write $P_{oo}=LL^{\mathsf T}$,
$P_{oo,+}=L_+L_+^{\mathsf T}$, and set
\begin{align}
 M&=L_+^{-1}AL,& N&=L_+^{-1}B,\nonumber\\
 Q&=\rho\Id-M^{\mathsf T}M\succ0.
\end{align}
Completing the square gives the sharp linear gain
\begin{equation}
 \gamma_b=\lambda_{\max}\!\left(
 N^{\mathsf T}N+(M^{\mathsf T}N)^{\mathsf T}Q^{-1}(M^{\mathsf T}N)
 \right).
 \label{eq:bias-schur}
\end{equation}
This gain does not bound the unknown remainder in
\eqref{eq:complement-map}. Every-prefix retention also needs all-direction
operator and supply bounds at each prefix. Positive endpoint margin alone
establishes neither condition.

\section{Conditional practical-stability statement}
After capture, let $t_j$ be consecutive informative-service superword
boundaries and $e_j$ the complete finite estimation-error coordinate. A
completion certificate must establish, on one retained domain,
\begin{align}
 m\norm{e_j}^2&\le V_j\le M\norm{e_j}^2,\nonumber\\
 V_{j+1}&\le\rho V_j+c_d\norm{d}_{[t_j,t_{j+1}]}^2,
 \qquad 0\le\rho<1,
 \label{eq:diss}
\end{align}
together with finite every-prefix bounds and forward retention under every
implemented event. If $\norm d\le D$, iteration gives
\begin{equation}
 \limsup_{j\to\infty}\norm{e_j}
 \le\sqrt{\frac{c_d}{m(1-\rho)}}D .
\end{equation}
Zero supply gives geometric decay at service boundaries. On the held-bias leg
$e_j$ is the complement coordinate and $d$ retains $Be_b+r$, because
\eqref{eq:held-bias-obstruction} rules out the full-state form. This
implication is conditional until capture, held-bias finite-error dissipation in
its restated form, the H18--A21 release, A21 dissipation, recurring service,
retention, and implementation arithmetic are certified.

\balance
\section{Declared constants}
Symbolic assumptions are separated from current numerical certification limits
in Table~\ref{tab:limits}. A declared value is a testable deployment
requirement, not proof that every assembled unit satisfies it.

\begin{table}[t]
\centering
\caption{Current deterministic certification limits}
\label{tab:limits}
\small
\begin{tabular}{lll}
\toprule
Quantity & Value & Unit\\
\midrule
$P_{\max}$ & 8.10 & m\\
$V_{\max}$ & 5.50 & m/s\\
$A_{\max}$ & 8.80 & m/s$^2$\\
$\Omega_{\max}$ & 0.610865 & rad/s\\
$P_{\rm AC}$ & 1100 & m\,s\\
$B_a$ & 0.225167 & m/s$^2$\\
$D_a$ & 0.001 & m/s$^3$\\
$B_g$ & 0.020 & rad/s\\
$D_g$ & $1.0\times10^{-5}$ & rad/s$^2$\\
$R_b$ (estimate) & 0.4 & m/s$^2$\\
$\tau_{b,e}$ (estimate) & 5000 & s\\
$T_M$ & 1.0 & s\\
$\mu_M$ & 1.0 & normalized\\
$\norm B$ & 20--75 & $\mu$T\\
$B_{\rm horiz,min}$ & 15 & $\mu$T\\
hard-iron residual & 5 & $\mu$T\\
magnetic residual & 2 & $\mu$T\\
\bottomrule
\end{tabular}
\end{table}

The accelerometer bias-rate value is a selected deterministic qualification
limit; assembled-history qualification remains open. Magnetic information
constants likewise remain to be certified uniformly from actually applied
shipping sensitivities.

\section{Evidence boundary}
Executable regressions verify the constant-displacement exclusion, sampled
predecessor bias rule, mode-dependent prediction identity, separate
correction/projection algebra, attempted-versus-applied magnetic distinction,
rank-deficient frequent magnetic service, unipotent obstruction, exactness of
the bias hold across a magnetically served window, and state inheritance across
Live and bias release. A one-second shipping superword, rooted thirty
seconds after Live, gives a complement covariance-metric ratio $0.999803$.
The all-direction prefix maximum is $1.001024$. Magnetic sensitivities and
transport are sampled before each applied correction; the measured information
minimum is $1.707892$, above the declared floor of one.

These are local incremental measurements, not a finite-error certificate.
The whitened fine/coarse sensitivity discrepancy is $0.014011$, whereas the
available operator-norm margin is $9.846\times10^{-5}$. The discrepancy is not
a rigorous uncertainty bound; the positive measured margin is not yet
certified against sensitivity error. High-precision evaluation cannot recover
precision lost in the shipping differences. Finite-error supply, recurring
service, all-time motion membership, capture and retention remain open.
Finite simulations do not establish these all-time properties.

\section{Conclusion}
The stability problem is posed as one same-history theorem with continuing
magnetic information. The formulation excludes permanent wave-coordinate
offsets without excluding quiet water, treats total residual bias independently
of estimator priors, and uses one prediction equation across held and active
accelerometer-bias modes. End-to-end regional practical stability remains
conditional on the open capture, dissipativity, retention, qualification, and
arithmetic certificates.

\section*{Acknowledgments and disclosure}
This work was prepared with supervised assistance from artificial intelligence
tools for mathematical drafting, code development, and editorial preparation.
The author is responsible for the conceptual development and final revisions,
declares no conflicts of interest, and received no funding for this research.
\end{document}
